\documentclass[12pt,a4paper]{article}

\usepackage{fontspec}
\usepackage{xeCJK}
\usepackage[a4paper,top=3cm,bottom=2.54cm,left=3.175cm,right=3.175cm,
            footskip=1.15cm]{geometry}
\usepackage{fancyhdr}
\usepackage{caption}
\usepackage{amsmath,amssymb}
\usepackage{array,tabularx,longtable}
\usepackage{enumitem}
\usepackage{tikz}
\usepackage[american]{circuitikz}
\usepackage{float}
\usepackage{xcolor}
\usepackage{listings}
\usetikzlibrary{arrows.meta,positioning,shapes.geometric,fit}

% 中文统一为宋体，西文和数字统一为 Times New Roman。
% 宋体用 macOS 自带的 Songti SC，Windows/Overleaf 上可换成 SimSun。
\setmainfont{Times New Roman}
\setmonofont{Times New Roman}
\setCJKmainfont[AutoFakeBold=2.5]{Songti SC}
\setCJKsansfont[AutoFakeBold=2.5]{Songti SC}
\setCJKmonofont{Songti SC}

% 三号 16bp（大标题）、四号 14bp（小标题）、小四 12bp（正文，行距固定 22bp）、
% 五号 10.5bp（图名表名）。
\newcommand{\xiaosi}{\fontsize{12bp}{22bp}\selectfont}
\newcommand{\wuhao}{\fontsize{10.5bp}{15bp}\selectfont}
\newcommand{\placeholder}{\textcolor{gray}{待实测}}

% 附录代码用等宽西文字体。附录列出的源文件全部是 ASCII，代码框里不会出现中文。
\newfontfamily\codefont{Courier New}

\lstdefinelanguage{Rust}{
  keywords={as,async,await,break,const,continue,crate,dyn,else,enum,extern,
    false,fn,for,if,impl,in,let,loop,match,mod,move,mut,pub,ref,return,self,
    Self,static,struct,super,trait,true,type,unsafe,use,where,while},
  morekeywords=[2]{bool,char,f32,f64,i8,i16,i32,i64,isize,str,u8,u16,u32,u64,
    usize,Option,Some,None,Result,Ok,Err},
  sensitive=true,
  morecomment=[l]{//},
  morecomment=[s]{/*}{*/},
  morestring=[b]{"},
}

\lstset{
  language=Rust,
  basicstyle=\codefont\fontsize{8bp}{9.8bp}\selectfont,
  keywordstyle=\bfseries,
  keywordstyle=[2]\itshape,
  commentstyle=\color{black!58},
  stringstyle=\color{black!75},
  numbers=left,
  numberstyle=\codefont\fontsize{6.5bp}{9.8bp}\selectfont\color{black!45},
  numbersep=7pt,
  firstnumber=1,
  breaklines=true,
  breakatwhitespace=false,
  breakindent=0pt,
  % 折行标记只能是不带颜色的 \mbox：postbreak 里放 \color 或 \textcolor，
  % listings 量这段材料的宽度时会读到空值，整个清单一行一个 Missing number。
  postbreak={\mbox{$\hookrightarrow$\space}},
  columns=fullflexible,
  keepspaces=true,
  showstringspaces=false,
  tabsize=4,
  frame=none,
  captionpos=t,
  xleftmargin=2.2em,
  aboveskip=8pt,
  belowskip=4pt,
}
\renewcommand{\lstlistingname}{代码}

\setlength{\parindent}{2em}
\setlength{\parskip}{0pt}
\setlength{\tabcolsep}{4pt}
\renewcommand{\arraystretch}{1.3}
\setlist{nosep}

\pagestyle{fancy}
\fancyhf{}
\renewcommand{\headrulewidth}{0pt}
\fancyfoot[R]{\wuhao\thepage}

% 图名、表名与正文同字体，字号五号。
\DeclareCaptionFont{wuhaocap}{\fontsize{10.5bp}{15bp}\selectfont}
\captionsetup{
  font=wuhaocap,
  labelfont=wuhaocap,
  labelsep=space,
  justification=centering,
  singlelinecheck=true,
  skip=3pt
}
\renewcommand{\figurename}{图}
\renewcommand{\tablename}{表}

\newcommand{\mainheading}[1]{%
  \par\vspace{6pt}\noindent{\fontsize{16bp}{22bp}\selectfont\bfseries #1}\par
  \vspace{3pt}}
\newcommand{\subheading}[1]{%
  \par\vspace{2pt}\noindent{\fontsize{14bp}{22bp}\selectfont\bfseries #1}\par}
\newcommand{\thirdheading}[1]{%
  \par\noindent{\xiaosi\bfseries #1}\par}

\tikzset{
  block/.style={draw,align=center,minimum height=8mm,minimum width=20mm},
  decision/.style={draw,diamond,aspect=2,align=center,inner sep=1pt},
  line/.style={-{Latex[length=2mm]},thin}
}

\begin{document}
\xiaosi

% 摘要单独成页，不显示页码；正文从第 1 页开始。
\thispagestyle{empty}
\begin{center}
  {\fontsize{16bp}{22bp}\selectfont\bfseries 摘要}
\end{center}

电流表由电流互感器取能，次级 270 匝并设抽头，其中 170 匝经 SS16 全桥整流、齐纳箝位
后向储能电容充电，再由低静态电流 LDO 稳压为 3.3\,V，不接任何外部电源。取电与传感
分时复用同一绕组：空闲期 PMOS、NMOS 双端关断，互感器输出全部用于储能；测量时两管
导通接入采样电阻，MSPM0G3507 以内部 2.5\,V 基准、4\,kHz 采样率经 DMA 采集 800 点，
用 Hann 加权真有效值和分档标定表换算出电流，约 260\,ms 后按序掉电。创新点是以一只
隔离肖特基把储能支路与采集支路解耦，使一次测量完全由电容独立供电；并以实测标定表
代替名义比例系数，吸收互感器非线性。读表器为 Android 应用，经
HC-42 蓝牙以请求-应答方式读数，支持人工与 120\,s 自动两种模式、越限报警和记录趋势。
0.55～2.32\,A 区间已建立 38 点标定表，表内拟合残差中位数 0.238\%，系统在无外部供电
下完成测量与无线上报。

\noindent\textbf{关键词：}无源电流表；MSPM0G3507；真有效值；分时取能；低功耗；BLE

\clearpage
\setcounter{page}{1}

\mainheading{一、\quad 方案比较与论证}

\subheading{1．取能与传感方案}

方案一：两组独立次级绕组，分别驱动取电电路和传感电路。两条通道互相隔离，调试直观，
但两组绕组同时占用磁芯窗口，绕制量大，匝数和耦合差异还会引入一致性误差。

方案二：一组带抽头的次级绕组，用受控 MOS 开关对取电和传感分时复用。空闲期断开传感
前端，互感器输出全部经全桥向储能电容充电；测量窗口内 PA26 拉低使 PMOS 导通、PA8
拉高使 NMOS 导通，传感前端接入，此时系统由储能电容独立供电。该方案不重复绕组，且传感
负载不在冷启动阶段占用取能功率——本题按启动电流计分，这一点直接影响得分。

本系统采用方案二。次级绕组总计 270 匝并设抽头，经不同抽头组合比较，以 170 匝有效
绕组接入整流桥。

\subheading{2．整流与稳压电源方案}

三个方案的前段相同：次级经全桥整流后向储能电容充电，齐纳二极管 $D_{\mathrm{Z}}$ 限制
储能节点电压。箝位不可省略，线路电流达到 2.2\,A 时次级开路电压远高于器件耐压。区别
只在储能节点之后靠什么得到 3.3\,V，电路如附录图 \ref{fig:power_options} 所示。

方案一：储能节点直接作为工作电压，不设稳压器。器件最少，静态电流只有齐纳的漏电流。
但箝位点即工作电压，储能电容只能充到 3.3\,V 附近，而电容储能与端电压平方成正比，同
容值下可存能量被压到最小。

方案二：箝位点取在远高于 3.3\,V 处，储能电容工作在该电压，再由低静态电流 LDO 降到
3.3\,V。储能电压提高之后同容值电容存的能量按平方增加，LDO 静态电流可做到微安级，
不需要电感，输出无开关纹波。代价是线性稳压搬运的是电荷而不是能量，压差部分在稳压器
上变成热。

方案三：稳压改用降压型开关变换器。变换器按功率守恒工作，同一段放电能交付
$\eta\cdot\frac{1}{2}C_{\mathrm{s}}(V_{\mathrm{H}}^{2}-V_{\mathrm{L}}^{2})$，高压段的
能量也取得出来。但变换器从储能电压刚开始建立时就在耗电：输入低于其最小启动电压期间
它并不工作，却仍有静态电流和开关节点漏电，把冷启动最紧张的那点取能功率消耗掉。


本系统采用方案二。冷启动阶段取能功率只有毫瓦量级，稳压器自身在这一阶段的消耗直接
决定电流表能不能启动，而启动电流与启动时间合计 40 分，权重高于稳压效率；LDO 在储能
电压尚未建立时几乎不取电，方案三则相反。测量窗口的能量由储能电容按式（4）承担，线性
稳压的效率损失不影响单次测量能否完成。

\subheading{3．电流有效值测量与量程方案}

方案一：外部偏置与外部运放。由独立基准把交流采样信号偏置到 ADC 输入范围中点，再经
外部运放放大。电路独立、调试直观，但需增加运放、偏置基准和增益切换器件，占板面积并
带来额外静态功耗。

方案二：使用 MSPM0G3507 片内 OPA 和 DAC。DAC 产生可调偏置，片内 OPA 完成放大，由
软件配置 1、2、4、8、16、32 倍增益并按采样幅值自动切换。外围器件少、静态功耗低，
偏置和增益可由程序协同调整。

本系统采用方案二，并用 ADC 采样序列计算真有效值，不使用整流滤波后测平均值再折算的
做法——后者只对正弦有效，而标定要对标的是标准表的真有效值读数。

\mainheading{二、\quad 参数计算、系统设计与实现}

\subheading{1．电路关键参数计算}

忽略励磁电流与损耗时，互感器原、副边安匝平衡，
\begin{equation}
  N_1 I_1=N_2 I_2 ,
\end{equation}
故副边电流为
\begin{equation}
  I_2=\frac{N_1}{N_2}I_1 .
\end{equation}
取 $N_1=5$、接入匝数 $N_2=170$，原边 2.0\,A 对应副边 58.8\,mA。增加匝数可提高低电流
下的感应电压，同时也增加绕组电阻和分布参数，故做成 270 匝带抽头、按实测选定接入匝数
的结构。

桥式整流每半周有两只二极管串联导通。设次级瞬时电压为 $u_2(t)$、单只二极管正向压降为
$V_{\mathrm{F}}$，则整流后瞬时电压为
\begin{equation}
  u_{\mathrm{r}}(t)=\max\left(|u_2(t)|-2V_{\mathrm{F}},\,0\right).
\end{equation}
输入电压越低，$2V_{\mathrm{F}}$ 占比越大，故全桥和隔离二极管一律选肖特基。

储能电容要独立支撑一次测量。设测量窗口内储能电压由 $V_{\mathrm{H}}$ 降到
$V_{\mathrm{L}}$，线性稳压交付负载的是电荷，故容量下限为
\begin{equation}
  V_{\mathrm{o}}\,C_{\mathrm{s}}\left(V_{\mathrm{H}}-V_{\mathrm{L}}\right)
  \geq E_{\mathrm{measure}} ,
\end{equation}
其中 $V_{\mathrm{o}}=3.3$\,V，$E_{\mathrm{measure}}$ 为一次约 260\,ms 窗口的能耗，
$V_{\mathrm{L}}$ 不得低于 LDO 压差加 3.3\,V。该式给出的是容量下限，实际取值以窗口内
$V_{\mathrm{H}}$、$V_{\mathrm{L}}$ 实测波形校核。

式（4）描述的是储能电压已经建立之后的一次测量。冷启动阶段储能电压自零升起，其间既不
具备完成一次测量的能量，稳压输出也尚未进入稳压区。固件据此设开机门限：上电后先测
$V_{\mathrm{DD}}$，达到 3.3\,V 才开始工作。测量借用 ADC0 的内部电源监测通道，该通道
不占引脚——选中它即断开外部输入并接到片内 $V_{\mathrm{DD}}/3$ 分压，与 ADC 共用同一
片内 2.5\,V 基准，故满量程为 $3\times2.5=7.5$\,V、分辨力 1.83\,mV，门限对应的码值为
\begin{equation}
  D_{\mathrm{th}}=\frac{V_{\mathrm{th}}}{3}\cdot\frac{2^{12}}{V_{\mathrm{ref}}}
  =\frac{3.3}{3}\cdot\frac{4096}{2.5}=1802 .
\end{equation}
2.5\,V 基准要求 $V_{\mathrm{DD}}\geq2.7$\,V，低于判决点，故判决本身落在基准的规格
范围内；供电更低时基准随供电下降，商值增大使读数偏高而非偏低，不会把未达门限的供电
误判为已达。

对增益为 $G$ 的同相偏置放大电路，输出为
\begin{equation}
  V_{\mathrm{o}}=G V_{\mathrm{i}}+(1-G)V_{\mathrm{DAC}} .
\end{equation}
若输入直流中心为 $V_{\mathrm{i,dc}}$，要求输出中心落在 ADC 中点 $V_{\mathrm{c}}$，则
$G>1$ 时
\begin{equation}
  V_{\mathrm{DAC}}=\frac{G V_{\mathrm{i,dc}}-V_{\mathrm{c}}}{G-1}.
\end{equation}
固件以 12 位 ADC 的 2047 码为目标中心，目标峰峰值利用率取满量程的 65\%，换档回差的
下、上限分别为 25\% 和 75\%。

TIMG6 以 4\,kHz 触发 ADC1，DMA 搬运采样。对 $N$ 点序列 $x[n]$ 施加 Hann 权 $w[n]$，
加权均值为
\begin{equation}
  \bar{x}_{w}=\frac{\sum_{n=0}^{N-1}w[n]x[n]}{\sum_{n=0}^{N-1}w[n]} ,
\end{equation}
去除直流后的有效值码宽为
\begin{equation}
 X_{\mathrm{rms}}=\sqrt{\frac{\sum_{n=0}^{N-1}w[n]
 \left(x[n]-\bar{x}_{w}\right)^2}{\sum_{n=0}^{N-1}w[n]}} .
\end{equation}
加窗是因为 200\,ms 记录未必覆盖整数个工频周期，Hann 权把截断处的权重压到零，使非整
周期截断不再直接进入结果。各量程分别保存标定点 $(x_k,I_k)$，当
$x_k\leq X_{\mathrm{rms}}\leq x_{k+1}$ 时按
\begin{equation}
 I=I_k+\frac{X_{\mathrm{rms}}-x_k}{x_{k+1}-x_k}(I_{k+1}-I_k)
\end{equation}
分段线性插值。采样点取在整流之后，波形是全波整流后的脉动量，其有效值与原边电流的
比例由标定表确定，不依赖对波形因数的解析假设，互感器非线性、采样电阻和链路增益的
误差一并被表吸收。

\subheading{2．硬件电路设计}

系统由电流互感器、SS16 全桥、齐纳箝位与储能、LDO、PMOS/NMOS 分时开关、片内模拟
前端、MSPM0G3507、OLED、受控供电的 HC-42 和 Android 读表器组成，框图如图
\ref{fig:system} 所示。电流表 A、B 结构相同，供能与测量链路互相独立，拆除任一表不
影响另一表。完整原理图见附录图 \ref{fig:sch}。

\begin{figure}[H]
\centering
\begin{tikzpicture}[node distance=4mm and 5mm,scale=0.58,transform shape]
  \node[block] (line) {36\,V 交流线路\\0.1～2.2\,A};
  \node[block,right=of line] (ct) {CT\\270 匝带抽头};
  \node[block,right=of ct] (bridge) {SS16 全桥\\箝位储能};
  \node[block,right=of bridge] (power) {LDO\\3.3\,V};
  \node[block,below=of bridge] (sw) {PMOS/NMOS 开关\\PA26/PA8};
  \node[block,right=of sw] (afe) {采样电阻与前端\\PA18/OPA1};
  \node[block,right=of afe] (mcu) {MSPM0G3507\\RMS/量程/地址};
  \node[block,below=of afe] (oled) {S2 按键\\OLED};
  \node[block,right=of mcu] (ble) {PB17 电源控制\\HC-42};
  \node[block,right=of ble] (app) {Android 读表器\\报警/存储/趋势};
  \draw[line] (line)--(ct);
  \draw[line] (ct)--(bridge);
  \draw[line] (bridge)--(power);
  \draw[line] (ct.south)|-(sw.west);
  \draw[line] (sw)--(afe);
  \draw[line] (afe)--(mcu);
  \draw[line] (power.south)|-(mcu.north);
  \draw[line] (mcu.south)|-(oled.east);
  \draw[line] (mcu)--(ble);
  \draw[line] (ble)--(app);
  \node[draw,dashed,fit=(ct)(bridge)(power)(sw)(afe)(mcu)(oled)(ble),
        label={[font=\wuhao]above:电流表 A（电流表 B 结构相同且独立）}] {};
\end{tikzpicture}
\caption{系统框图}
\label{fig:system}
\end{figure}

取能支路以 170 匝有效绕组经四只 SS16 组成的全桥整流，正、负半周均向储能电容充电。
整流输出记为节点 $A$，与储能电容之间另串一只肖特基 $D_5$：采集支路接入时把 $A$ 拉低，
若无 $D_5$，冷启动积攒的能量会经采集支路反向放掉；有了 $D_5$ 能量只能单向流入电容，
测量期间由电容独立维持系统。储能电容并接串了限流电阻的按钮，即题目要求的残余电量
放电回路。

采集支路自节点 $A$ 取出，由 PMOS 高边和 NMOS 低边双端把关，采样电阻 $R_{\mathrm{s}}$
串在两者之间，PMOS 漏极与 $R_{\mathrm{s}}$ 的交点送 PA18。两管同时导通时支路成回路，
$R_{\mathrm{s}}$ 上的压降随次级电流变化；任一只关断支路即断开，窗口之外不消耗取能
功率。PA26 经 10\,k$\Omega$ 上拉至 $A$，PA8 经 10\,k$\Omega$ 下拉至地：上电
到 MCU 接管 GPIO 之前引脚是输入而非驱动输出，这两只电阻使支路在该窗口内保持断开，而
这段窗口恰好是储能电压最弱、最经不起额外负载的时候。分时控制状态见附录表
\ref{tab:afe_state}。

模拟前端全部在片内：OPA1 作同相可编程增益放大，DAC12 给出偏置，ADC 和 DAC 的基准取
片内 2.5\,V 基准而不是供电电压，使供电纹波和 LDO 输出误差不进入读数。引脚分配见附录
表 \ref{tab:pins}。

\subheading{3．软件设计}

程序上电后先把 PA26/PA8 置为取电空闲态，PB17 拉低使蓝牙断电，OLED 不发送初始化命令，
随后停下来等电源：每 20\,ms 用 ADC0 读一次 $V_{\mathrm{DD}}$，不到 3.3\,V 就不往下
走。这一步之前除电源指示 LED 外没有任何模拟模块使能、蓝牙也未上电，储能电容的充电
不被负载打断；基准和 ADC0 本身也是读完即掉电，仅占 1.5\% 的占空比。等待不设超时——
供电达不到 3.3\,V 时，稳压输出尚未进入稳压区，OLED 与蓝牙模块都不在额定供电条件下，
停在这里等待比出一个不可信的读数、或让欠压的外设动作更合理，此时电源指示 LED 仍在
闪烁，与死机可区分。

电压到位后不等请求先完成一次测量——读表器一连上就有数可读。这次测量结束后才把 PB17
置高、等待 350\,ms 并初始化 UART2：模块的启动电流是本固件中最大的单笔负载，不应落在
供电刚建立的时刻。

此后主循环等待三件事之一：S2 下降沿、蓝牙收到 \texttt{MEAS}、2\,s 心跳定时到期。前两
者完全等价，走同一条测量路径；心跳只发一帧保活并继续等待。测量期间到达的请求两者
都丢弃而不排队，测量结束时清空串口积压字节，使下一次等待从静默开始。

一次测量的时序是：切换为传感态并等待 20\,ms 建立，采 160 点/40\,ms 探测帧，按峰峰值
和直流中心选定增益与 DAC 偏置（最多 1\,ms 建立），再采 800 点/200\,ms 正式帧；然后按
TIMG6/ADC1、OPA1/DAC12、基准、外部 MOS 开关的顺序关断，最后才计算有效值并查表。
关断顺序不可颠倒：先断外部开关会让最后一段采样落在正在塌陷的输入上。程序流程见附录
图 \ref{fig:flow}。

读数带质量标志：基准未就绪、输入越出转换窗口、采集帧未完成时，该次读数不可信，固件
只发调试帧并给出原因，不发读数帧，避免把不可信的数值以正常状态送出。

\subheading{4．通信协议设计}

链路为 HC-42 透传串口（FFE0 服务、FFE1 通知特征），速率 115200\,bit/s。透传链路不提供
分包与长度字段，故采用面向行的 ASCII 协议，以换行作自同步定界符：接收端在任意位置
接入或丢失字节后，都能在下一个换行重新对齐。题目 2(2) 规定两种模式下均不得操作电流
表，因此读数由读表器请求产生，电流表不主动上报测量值。三种帧如下。

\noindent（1）读表器 $\to$ 电流表，请求帧：

\centerline{\texttt{MEAS} \quad 或 \quad \texttt{MEAS,ADDR=dd}}

前者为广播，后者仅编码开关为 \texttt{dd} 的电流表应答。命令带地址而非仅用广播，是
因为一台电流表在不同时刻由编码开关赋予不同地址，仅凭连接无法确定对端身份；带地址后
请求与应答一一对应，不会把某地址的应答记到另一地址名下。

\noindent（2）电流表 $\to$ 读表器，读数帧与调试帧：

\centerline{\texttt{METER\_TEST,ADDR=dd,CURRENT\_MA=iiii,STATUS=ssss}}
\centerline{\texttt{METER\_CAL,ADDR=dd,RMS=257.83,GAIN=1,FLAG=OK,MEAN=501,PP=856}}

\texttt{CURRENT\_MA} 为有效值，单位 mA；\texttt{STATUS} 取 \texttt{NORMAL}、
\texttt{LOW}、\texttt{HIGH}，按 0.2\,A 与 2.0\,A 判定，没有 \texttt{OFFLINE}——离线是
读表器的判断，电流表无法报告自己不在。调试帧与读数帧同次产生，\texttt{RMS} 和
\texttt{GAIN} 用于建立标定表，\texttt{MEAN}、\texttt{PP} 是探测帧测得的输入直流与峰峰
值，\texttt{FLAG} 是该次读数的质量标志。

\noindent（3）电流表 $\to$ 读表器，保活帧：

\centerline{\texttt{IMHERE,ADDR=dd}}

空闲时每 2\,s 一次，不含读数，也不唤醒任何模拟电路。作用有二：编码开关改变时旧地址的
心跳停止、新地址出现，读表器据此更新在场地址；负载断开后电流表随之失电，心跳消失即
为题目 2(3) 的离线判据，比等待一次请求超时更快。请求超时只说明本次未读到，不据此
判定离线。

\subheading{5．其他重要模块设计}

OLED 上电不初始化，一次测量结束后显示 1\,s 即熄屏，对比度设为最低档，空闲期不刷屏。
蓝牙模块不是仅延后 UART，而是由 PB17 控制外部 3.3\,V 负载开关，低电平关断、高电平
导通；断电期间 UART 引脚保持未初始化，避免经模块 RXD 保护二极管反向灌电。地址由
PB0、PB6～PB8 四位编码开关给出，低有效，范围 00～15，每次心跳都重读，使运行中拨动
开关立即生效。

Android 读表器扫描 8\,s、间歇 12\,s，连接名称符合 HC-42、METER、AMMETER 等前缀的
设备。基本模式为一对一手动读取；自动模式一键启动，每 120\,s 对在场地址依次请求，同一
连接上只允许一个未完成请求。读数写入 SQLite，容量 1000 条，历史与趋势均超过要求的
30 条；离线与请求超时都不写入记录——记录用于保存读数，缺席不是读数。低于 0.2\,A 与
高于 2.0\,A 的每一次读数都报警，不以状态跳变为条件：读数稀疏，同一状态的相邻两次读数
是两个独立事件。

\mainheading{三、\quad 测试方案及测试结果分析}

\subheading{1．测试仪器与测试条件}

测试在室温下进行，由 220\,V/36\,V、60\,VA 单相隔离变压器供电，回路串入标定电流表和
自制负载，负载视在功率不超过 50\,VA、电流 0～2.2\,A 连续可调。仪器见表
\ref{tab:instruments}。

\begin{table}[H]
\centering
\caption{测试仪器列表}
\label{tab:instruments}
\wuhao
\begin{tabularx}{\textwidth}{|>{\centering\arraybackslash}p{8mm}|
>{\centering\arraybackslash}p{28mm}|>{\centering\arraybackslash}p{30mm}|
>{\centering\arraybackslash}p{8mm}|X|}
\hline
序号 & 仪器类别 & 仪器型号 & 数量 & 用途\\ \hline
1 & 交流电源 & 熙硕 STG-500W & 1 & 提供可调交流测试电源\\ \hline
2 & 五位半万用表 & SIGLENT SDM3055X-E & 1 & 标准电流与电压比对\\ \hline
3 & 数字示波器 & SIGLENT SDS2102X PLUS & 1 & 储能、采样与启动时序\\ \hline
4 & Android 手机 & 通用 & 1 & 无线读取、报警与存储测试\\ \hline
5 & 秒表 & 通用 & 1 & 记录启动与通信时间\\ \hline
\end{tabularx}
\end{table}

\subheading{2．测试方法}

（1）精度测试：保持线路电压 36\,V，依次设置 0.10、0.20、0.50、1.00、1.50、2.00\,A，
每点稳定后同时记录标准表与被测表读数，每点重复不少于 5 次，A、B 两表分别测试；再
拆除 B 表重复测 A，比较误差变化。

（2）标定与复现性测试：以脚本同步采集 BLE 帧和 SDM3055X-E 读数，负载从低到高再从高
到低连续扫描一遍，逐帧记录标准电流、RMS 码宽、量程和质量标志，用于建表并统计帧间
散布与来回两趟的系统差。

（3）取能与启动测试：储能电容经放电按钮充分放电后，从小到大增加负载电流，记录能够
稳定完成一次测量并发帧的最低电流，即启动电流；在该电流下用示波器捕获储能电压建立、
首次有效值生成和首帧到达的时刻，得到启动时间。

（4）分时时序测试：同时观察 PA26、PA8、储能电压和 3.3\,V，确认空闲态为 PA26=1、
PA8=0，测量态为 PA26=0、PA8=1，并记录 260\,ms 传感窗口内的储能电压跌落。

（5）显示与无线测试：确认上电阶段 OLED 无驱动，单击 S2 后显示 1\,s 自动熄屏；手机与
电流表相距 5\,m 检查人工读取和 120\,s 自动模式；设置低限、正常、超限并断开一只表检查
报警；写入 30 条以上记录后重启应用，检查历史与趋势。

\subheading{3．测试结果}

供电绕组：85 匝单段绕组感应电压和调节余量不足；改为 270 匝总绕组、每 50 匝预留抽头
后可组合不同有效匝数，经对比最终以 170 匝有效绕组供电。

标定与复现性：2026 年 7 月 31 日的连续扫描共 196 帧，其中 192 帧质量标志为
\texttt{OK}，标准电流覆盖 0.5503～2.3218\,A。归并为 38 个标定点写入 X1 直采档的标定
表，覆盖 0.5527～2.2633\,A。表内拟合残差中位数 0.238\%、90 分位 0.714\%；同一电流点
的帧间散布 1$\sigma$ 为 0.45\%；下降趟在同一电流上比上升趟低 0.59\%，把两趟合并入表
后，上升点残差 $+0.17\%$、下降点残差 $-0.24\%$。剩余 4 帧标志为 \texttt{OVER}：
2.33\,A 时探测帧测得波形峰值已达 4107 码，越过 4095 上限。

精度记录见表 \ref{tab:accuracy}，按标准表比对结果逐格填写，不填入推测数据。启动电流、
启动时间、5\,m 通信距离和外形尺寸同样以实测填写。

\begin{table}[H]
\centering
\caption{电流测量精度记录表}
\label{tab:accuracy}
\wuhao
\begin{tabular}{|c|c|c|c|c|c|c|}
\hline
标准值/A & 0.10 & 0.20 & 0.50 & 1.00 & 1.50 & 2.00\\ \hline
A 表示值/A & \placeholder & \placeholder & \placeholder & \placeholder & \placeholder & \placeholder\\ \hline
A 表误差/\% & \placeholder & \placeholder & \placeholder & \placeholder & \placeholder & \placeholder\\ \hline
B 表示值/A & \placeholder & \placeholder & \placeholder & \placeholder & \placeholder & \placeholder\\ \hline
B 表误差/\% & \placeholder & \placeholder & \placeholder & \placeholder & \placeholder & \placeholder\\ \hline
启动电流/A & \multicolumn{3}{c|}{\placeholder} & 启动时间/s & \multicolumn{2}{c|}{\placeholder}\\ \hline
\end{tabular}
\end{table}

\subheading{4．结果分析}

误差主要来自四处。其一是帧间散布 0.45\%：一次读数只用一帧 200\,ms 记录，工频与采样
窗口不同步、负载本身的波动都落在里面；Hann 加权已经压掉了非整周期截断的边界项，进一
步降低要靠多帧取中位数，但那是成倍的采样能耗，与启动电流指标冲突，故未采用。其二是
上下行 0.59\% 的系统差，其中一部分是比对时序造成的：标准表读数比电流表的采样窗口滞后
约 300\,ms，扫描时上升趟与下降趟的偏移方向相反；另一部分来自磁芯磁滞。把两趟合并入
同一张表使残差从单向的 0.59\% 降到双向的 $+0.17\%/-0.24\%$，是当前处理方式。

其三是标定覆盖：0.55\,A 以下尚无标定点，而该段信号幅度按比例下降，帧间散布的相对值
随之放大，是全量程 0.5\% 的主要风险；改进措施是在 0.1～0.55\,A 补点，并把采样点改到
隔直之后、用固定偏置送入 OPA，使直流中心与信号幅度解耦。当前前端从整流节点取样，
直流与交流同比例变化，PGA 会同时放大两者，因此实测中量程始终停在 X1 档——输入自身
的摆幅在每个电流下都占直流留出余量的 76\%～89\%，没有哪一档增益放得进去。其四是上端
余量：2.33\,A 时峰值已触及 4095 上限，2\,A 以上的越限区只有约 15\% 宽，改进措施是减小
采样电阻或抬高偏置留出的上行余量。

\mainheading{四、\quad 总结}

本系统以电流互感器取能，全桥整流、齐纳箝位、电容储能后由低静态电流 LDO 稳压供电，
取电与传感用一对 MOS 开关分时复用同一绕组，并以隔离肖特基保证测量窗口内由电容独立
供电。测量用片内 OPA、DAC 和 ADC，4\,kHz 采样 800 点，Hann 加权真有效值配合分段线性
标定表换算，读数经 HC-42 以请求-应答方式送给 Android 读表器。方案可行：0.55～2.32\,A
已建标定表，表内残差中位数 0.238\%，分时供电与无线上报均已跑通。优点是空闲期无传感
负载、模拟模块全部掉电，有利于启动电流指标，标定表把互感器非线性一并吸收；不足是低端
尚未标定、量程档位在当前前端上无法启用，需按上节措施调整采样点位置后重新标定。

\clearpage
\mainheading{附录}

\subheading{1．硬件电路原理图}

图 \ref{fig:sch} 中同名网络（$A$、3.3\,V、GND 以及各引脚名）相连。

\begin{figure}[H]
\centering
\begin{circuitikz}[line cap=round,line join=round,font=\wuhao,scale=0.76,
    transform shape,>={Latex[length=1.8mm]},
    blk/.style={draw,align=center,inner sep=2.5pt,minimum height=7mm}]
  \ctikzset{bipoles/length=0.8cm}
  % 桥臂是斜的，标签默认跟着元件转；straight 让它们保持水平。
  \ctikzset{label/align=straight}
  \def\Ro{1.32}\def\Ri{0.84}\def\Rm{1.08}\def\Ra{0.40}
  % 绕线是绕在环截面上的螺线：设绕过截面的相位为 p，正视投影半径为
  % \Rm+\Ra*cos(p)，p 在 0~180 之间位于环前面。先整条画出，再用不透明的
  % 磁环盖住背面那半，最后把前面那半重画——遮挡由磁环自身完成。
  \newcommand{\coilpath}[3]{%
    \pgfmathsetmacro{\kk}{#3*360/(#2-#1)}%
    \draw[thick,domain=#1:#2,samples=460,variable=\t,smooth]
      plot ({(\Rm+\Ra*cos(\kk*(\t-#1)))*cos(\t)},
            {(\Rm+\Ra*cos(\kk*(\t-#1)))*sin(\t)});}
  \newcommand{\coilfront}[3]{%
    \pgfmathsetmacro{\kk}{#3*360/(#2-#1)}%
    \pgfmathsetmacro{\jmax}{#3-1}%
    \foreach \j in {0,...,\jmax}{%
      \pgfmathsetmacro{\ta}{#1+\j*360/\kk}%
      \pgfmathsetmacro{\tb}{#1+(\j*360+180)/\kk}%
      \draw[thick,domain=\ta:\tb,samples=40,variable=\t,smooth]
        plot ({(\Rm+\Ra*cos(\kk*(\t-#1)))*cos(\t)},
              {(\Rm+\Ra*cos(\kk*(\t-#1)))*sin(\t)});}}

  % ---- 电流互感器 ----
  \coilpath{130}{230}{5}
  \coilpath{-70}{70}{9}
  \fill[black!10,even odd rule] (0,0) circle (\Ro) (0,0) circle (\Ri);
  \draw[thick] (0,0) circle (\Ro);
  \draw[thick] (0,0) circle (\Ri);
  \coilfront{130}{230}{5}
  \coilfront{-70}{70}{9}
  \draw[thick] (130:{\Rm+\Ra}) -- (-2.45,1.10);
  \draw[thick] (230:{\Rm+\Ra}) -- (-2.45,-1.10);
  \draw[thick,->] (-3.15,1.10) -- (-2.45,1.10);
  \node[above] at (-2.85,1.14) {$i_1$};
  \node[left] at (-2.50,-1.10) {被测线路};
  \node[below,align=center] at (0,-2.00)
    {取能/传感铁芯\\$N_1=5$，$N_2=270$（170 匝抽头接入）};
  \draw[thick,rounded corners=2pt] (70:{\Rm+\Ra}) -- (2.15,1.32) -- (2.15,0.55)
    -- (3.05,0.55);
  \draw[thick,rounded corners=2pt] (-70:{\Rm+\Ra}) -- (2.15,-1.32)
    -- (2.15,-0.55) -- (3.05,-0.55);

  % ---- SS16 全桥：n/s 为交流端，w 为负、e 为正 ----
  \coordinate (bn) at (5.40, 1.20);
  \coordinate (bs) at (5.40,-1.20);
  \coordinate (bw) at (4.20, 0);
  \coordinate (be) at (6.60, 0);
  \draw (bn) to[sD,l=$D_1$] (be);
  \draw (bs) to[sD,l_=$D_2$] (be);
  \draw (bw) to[sD,l=$D_3$] (bn);
  \draw (bw) to[sD,l_=$D_4$] (bs);
  \draw (3.05,0.55) -- (3.70,0.55) -- (3.70,1.90) -- (5.40,1.90) -- (bn);
  \draw (3.05,-0.55) -- (3.70,-0.55) -- (3.70,-1.90) -- (5.40,-1.90) -- (bs);
  \draw (bw) -- (3.20,0) -- (3.20,-3.05);

  % ---- 隔离、箝位、储能、放电与稳压 ----
  \draw (be) -- (7.55,0);
  \node[circ] at (7.55,0) {};
  \node[above] at (7.55,0.10) {$A$};
  \draw (7.55,0) to[sD*,l=$D_5$] (9.25,0);
  \draw (9.25,0) -- (12.05,0);
  \draw (9.45,0) to[C] (9.45,-3.05);
  \node[left] at (9.32,-1.55) {$C_{\mathrm{s}}$};
  \draw (10.40,0) to[R,l_=$R_{\mathrm{d}}$] (10.40,-1.60);
  \draw (10.40,-1.60) to[push button,l_=放电] (10.40,-3.05);
  \draw (11.85,-3.05) to[zD] (11.85,0);
  \node[right] at (11.98,-1.55) {$D_{\mathrm{Z}}$};
  \node[circ] at (9.45,0) {};
  \node[circ] at (10.40,0) {};
  \node[circ] at (11.85,0) {};
  \node[blk,minimum width=1.3cm] (ldo) at (12.85,0) {LDO};
  \draw (12.05,0) -- (ldo.west);
  \draw (ldo.east) -- (14.55,0);
  \draw (ldo.south) -- (12.85,-3.05);
  \node[circ] at (12.85,-3.05) {};
  \draw (14.10,0) to[C] (14.10,-3.05);
  \node[circ] at (14.10,0) {};
  \node[above] at (14.45,0.10) {3.3\,V};
  \draw (3.20,-3.05) -- (14.10,-3.05);
  \node[ground] at (8.30,-3.05) {};

  % ---- 采集支路 ----
  \node[above] at (0.55,-4.35) {$A$};
  \draw (0.55,-4.40) -- (0.55,-5.05);
  \node[pmos] (QP) at (0.55,-5.55) {};
  \draw (QP.S) -- (0.55,-5.05);
  \node[right=0.5mm of QP] {$Q_{\mathrm{P}}$};
  \draw (QP.D) -- (0.55,-6.45);
  \draw (0.55,-6.45) to[R,l_=$R_{\mathrm{s}}$] (0.55,-7.75);
  \node[nmos] (QN) at (0.55,-8.25) {};
  \draw (QN.D) -- (0.55,-7.75);
  \draw (QN.S) -- (0.55,-9.15) node[ground] {};
  \node[right=0.5mm of QN] {$Q_{\mathrm{N}}$};
  \draw (QP.G) -- (-2.35,-5.55) node[left] {PA26};
  \draw (-1.35,-5.55) to[R,l=$10\,\mathrm{k}\Omega$] (-1.35,-4.40) -- (0.55,-4.40);
  \draw (QN.G) -- (-2.35,-8.25) node[left] {PA8};
  \draw (-1.35,-8.25) to[R,l_=$10\,\mathrm{k}\Omega$] (-1.35,-9.15)
    -- (0.55,-9.15);

  % ---- 控制器与外设 ----
  \node[draw,minimum width=3.6cm,minimum height=4.3cm,align=center]
    (mcu) at (5.65,-6.90) {MSPM0G3507\\[2pt] OPA1/DAC12\\ ADC1/TIMG6/DMA};
  % 采样点直接进 PA18；两个栅极控制脚用同名网络标注，免得连线穿过采集支路。
  \node[circ] at (0.55,-6.45) {};
  \draw (0.55,-6.45) -- (3.85,-6.45);
  \node[above] at (2.20,-6.42) {PA18};
  \draw (3.35,-5.55) -- (3.85,-5.55);
  \node[left] at (3.38,-5.55) {PA26};
  \draw (3.35,-7.55) -- (3.85,-7.55);
  \node[left] at (3.38,-7.55) {PA8};
  \draw[->] (5.65,-4.15) -- (5.65,-4.75);
  \node[above] at (5.65,-4.10) {3.3\,V};
  \draw (5.65,-9.05) -- (5.65,-9.55) to[push button] (5.65,-10.35)
    node[ground] {};
  \node[left] at (5.45,-9.95) {S2（PB21）};

  \node[blk,minimum width=2.6cm] (oled) at (11.30,-5.20) {OLED\\SSD1306};
  \node[blk,minimum width=2.6cm] (bt)   at (11.30,-6.40) {HC-42\\BLE 模块};
  \node[blk,minimum width=2.6cm] (lsw)  at (11.30,-7.60) {3.3\,V 负载开关};
  \node[blk,minimum width=2.6cm] (dip)  at (11.30,-8.80) {4 位编码开关};
  \draw[->] (7.45,-5.20) -- (oled.west);
  \node[above] at (8.72,-5.16) {PB2/PB3};
  \draw[->] (7.45,-6.40) -- (bt.west);
  \node[above] at (8.72,-6.36) {PB15/PB16};
  \draw[->] (7.45,-7.60) -- (lsw.west);
  \node[above] at (8.72,-7.56) {PB17};
  \draw[->] (dip.west) -- (7.45,-8.80);
  \node[above] at (8.72,-8.76) {PB0/PB6$\sim$PB8};
  \draw[->] (lsw.north) -- (bt.south);
  \draw[->] (13.70,-7.60) -- (lsw.east);
  \node[above] at (13.25,-7.56) {3.3\,V};
\end{circuitikz}
\caption{电流表硬件原理图}
\label{fig:sch}
\end{figure}

\begin{figure}[H]
\centering
\begin{circuitikz}[line cap=round,line join=round,font=\wuhao,scale=0.88,
    transform shape,>={Latex[length=1.8mm]}]
  % 元件默认长度 1.4 cm，占满整条支路会显得笨重，这里压到 0.75 cm。
  \ctikzset{bipoles/length=0.75cm}
  % 三个方案的前段一样，用 \foreach 画一遍：整流输出 A 充储能电容，齐纳箝位。
  \foreach \x in {0,5.2,10.4}{
    \begin{scope}[xshift=\x cm]
      \draw[->] (-0.55,0) -- (0.70,0);
      \node[above] at (-0.15,0.02) {$A$};
      \draw (0.70,0) -- (1.85,0);
      \draw (0.85,0) to[C] (0.85,-1.50);
      \node[left] at (0.72,-0.75) {$C_{\mathrm{s}}$};
      % 由下往上画，齐纳的阴极才在上：箝位靠的是反向击穿，画反了就成了一只
      % 把储能节点直接短到地的正向二极管。
      \draw (1.85,-1.50) to[zD] (1.85,0);
      \node[right] at (1.99,-0.75) {$D_{\mathrm{Z}}$};
      \node[circ] at (0.85,0) {};
      \node[circ] at (1.85,0) {};
      \draw (0.85,-1.50) -- (1.85,-1.50);
      \node[ground] at (1.35,-1.50) {};
    \end{scope}}

  \begin{scope}[xshift=0cm]
    \draw (1.85,0) -- (3.60,0);
    \node[above] at (3.55,0.04) {3.3\,V};
    \node at (1.90,-2.15) {(a) 箝位直供};
  \end{scope}

  \begin{scope}[xshift=5.2cm]
    \node[draw,minimum width=1.15cm,minimum height=0.78cm,inner sep=1pt]
      (ldo) at (2.80,0) {LDO};
    \draw (1.85,0) -- (ldo.west);
    \draw (ldo.east) -- (4.00,0);
    \draw (3.70,0) to[C] (3.70,-1.50);
    \node[circ] at (3.70,0) {};
    \draw (1.85,-1.50) -- (3.70,-1.50);
    \node[above] at (3.95,0.04) {3.3\,V};
    \node at (1.90,-2.15) {(b) 低静态电流 LDO};
  \end{scope}

  \begin{scope}[xshift=10.4cm]
    \node[draw,minimum width=1.30cm,minimum height=0.78cm,inner sep=1pt]
      (sw) at (2.75,0) {DC-DC};
    \draw (1.85,0) -- (sw.west);
    \draw (sw.east) to[L] (4.30,0);
    \draw (4.30,0) -- (4.60,0);
    \draw (4.30,0) to[C] (4.30,-1.50);
    \node[circ] at (4.30,0) {};
    \draw (1.85,-1.50) -- (4.30,-1.50);
    \node[above] at (4.55,0.04) {3.3\,V};
    \node at (2.20,-2.15) {(c) 降压 DC-DC};
  \end{scope}
\end{circuitikz}
\caption{三种稳压方案}
\label{fig:power_options}
\end{figure}

\begin{table}[H]
\centering
\caption{取电与传感分时控制状态}
\label{tab:afe_state}
\wuhao
\begin{tabularx}{0.95\textwidth}{|>{\centering\arraybackslash}p{20mm}|
>{\centering\arraybackslash}p{14mm}|>{\centering\arraybackslash}p{14mm}|
>{\centering\arraybackslash}p{24mm}|X|}
\hline
工作状态 & PA26 & PA8 & MOS 状态 & 能量与信号路径\\ \hline
空闲/取电 & 高 & 低 & 均关断 & 传感前端隔离，互感器输出全部用于整流储能\\ \hline
测量/传感 & 低 & 高 & 均导通 & 传感前端接入，储能电容独立维持系统供电\\ \hline
\end{tabularx}
\end{table}

\clearpage
\subheading{2．程序流程图}

\begin{figure}[H]
\centering
\begin{tikzpicture}[scale=0.80,transform shape,font=\wuhao,
    lb/.style={draw,align=center,text width=4.2cm,minimum height=1.0cm,
               inner sep=2.5pt},
    rb/.style={draw,align=center,text width=4.7cm,minimum height=1.0cm,
               inner sep=2.5pt}]
  % 主流程在左，一次测量的子流程在右。两列之间不连线：测量由主流程调用两次
  % （上电自测和收到触发），画成子流程比拉两条回线清楚。
  \node[lb] (b1) at (0, 0)     {上电：PA26=1，PA8=0\\ PB17=0，OLED 不初始化};
  \node[decision] (d0) at (0,-1.90) {ADC0 读 $V_{\mathrm{DD}}$\\ $\geq3.3$\,V？};
  \node[lb] (b2) at (0,-3.85)  {上电自测一次\\ 执行测量子流程};
  \node[lb] (b3) at (0,-5.50)  {PB17=1，等 350\,ms\\ 初始化 UART2};
  \node[lb] (b4) at (0,-7.15)  {发读数帧与调试帧\\ OLED 显示 1\,s 后熄屏};
  \node[lb] (b5) at (0,-8.80)  {清空积压请求，等待触发};
  \node[decision] (d6) at (0,-10.60) {触发来源？};
  \node[lb,text width=2.8cm] (b7) at (-5.20,-10.60)
    {发 \texttt{IMHERE} 保活帧};
  \node[lb] (b8) at (0,-12.35) {执行测量子流程};

  \node[rb,rounded corners=8pt] (m0) at (9.00, 0)    {子流程：一次测量};
  \node[rb] (m1) at (9.00,-1.55)  {PA26=0，PA8=1，等 20\,ms 建立};
  \node[rb] (m2) at (9.00,-3.20)  {初始化基准与 DAC、OPA1、ADC1/TIMG6};
  \node[rb] (m3) at (9.00,-4.85)  {160 点/40\,ms 探测\\ 选增益与 DAC 偏置};
  \node[rb] (m4) at (9.00,-6.50)  {800 点/200\,ms DMA 正式采样};
  \node[rb] (m5) at (9.00,-8.15)  {关 ADC/TIMG6、OPA/DAC、基准\\ PA26=1，PA8=0};
  \node[rb] (m6) at (9.00,-9.80)  {Hann 加权真有效值\\ 分档标定表插值};
  \node[rb,rounded corners=8pt] (m7) at (9.00,-11.30) {返回读数};

  \draw[line] (b1)--(d0);
  \draw[line] (d0.south)--node[right]{是}(b2.north);
  % 否分支绕回判决自身：等待没有超时，供电不到 3.3 V 就一直停在这一步。
  \draw[line] (d0.west)--node[above]{否，等 20\,ms}(-4.60,-1.90)--(-4.60,-0.95)
              --(0,-0.95)--(d0.north);
  \draw[line] (b2)--(b3);
  \draw[line] (b3)--(b4);
  \draw[line] (b4)--(b5);
  \draw[line] (b5)--(d6);
  \draw[line] (d6.west)--node[above]{2\,s 心跳}(b7.east);
  \draw[line] (b7.north)--(-5.20,-8.80)--(b5.west);
  \draw[line] (d6.south)--node[right]{S2 或 \texttt{MEAS}}(b8.north);
  \draw[line] (b8.east)--(3.90,-12.35)--(3.90,-7.15)--(b4.east);
  \draw[line] (m0)--(m1);
  \draw[line] (m1)--(m2);
  \draw[line] (m2)--(m3);
  \draw[line] (m3)--(m4);
  \draw[line] (m4)--(m5);
  \draw[line] (m5)--(m6);
  \draw[line] (m6)--(m7);
\end{tikzpicture}
\caption{程序流程图}
\label{fig:flow}
\end{figure}

\subheading{3．控制器接口}

\begin{table}[H]
\centering
\caption{MSPM0G3507 主要接口分配}
\label{tab:pins}
\wuhao
\begin{tabularx}{\textwidth}{|>{\centering\arraybackslash}p{25mm}|
>{\centering\arraybackslash}p{30mm}|X|}
\hline
引脚 & 功能 & 说明\\ \hline
PA26 & PMOS 高边控制 & 低电平导通；10\,k$\Omega$ 上拉，上电默认关断\\ \hline
PA8 & NMOS 低边控制 & 高电平导通；10\,k$\Omega$ 下拉，上电默认关断\\ \hline
PA18 & ADC1 输入 & 采样电阻上的电流采样通道\\ \hline
PA16 & OPA1 输出 & 可编程增益后的模拟通道\\ \hline
PB2、PB3 & 软件 I\textsuperscript{2}C & OLED 的 SCL、SDA\\ \hline
PB21 & S2 按键 & 下降沿触发一次测量，与蓝牙 \texttt{MEAS} 等价\\ \hline
PB17 & 蓝牙电源使能 & 高电平导通外部负载开关，低电平关闭\\ \hline
PB15 & UART2 TX & HC-42 串口发送，115200\,bit/s\\ \hline
PB16 & UART2 RX & 接收读表器的 \texttt{MEAS} 请求\\ \hline
PB0、PB6～PB8 & 地址输入 & 4 位低有效地址，范围 00～15\\ \hline
\end{tabularx}
\end{table}

\clearpage
\subheading{4．主要程序代码}

电流表固件用 Rust 编写，编译目标为 \texttt{thumbv6m-none-eabi}，异步执行器为
embassy。以下按一次测量经过的顺序列出四个源文件：\texttt{main.rs} 决定何时测量，
\texttt{meter.rs} 执行一次测量的完整时序，\texttt{dsp.rs} 把采样序列算成有效值码宽，
\texttt{link.rs} 管地址、帧格式和收发。其余模块不在此列出：\texttt{sampler.rs} 配置
TIMG6/ADC1/DMA，\texttt{range.rs} 配置 OPA1 增益与 DAC 偏置并实现换档判据，
\texttt{cal.rs} 存分档标定表，\texttt{vref.rs} 开关内部基准，\texttt{supply.rs} 用
ADC0 的内部通道读供电电压，\texttt{oled.rs} 与 \texttt{ui.rs} 驱动显示。行号即源文件
行号。

\lstinputlisting[caption={\texttt{src/main.rs}：触发与调度}]{../src/main.rs}

\clearpage
\lstinputlisting[caption={\texttt{src/meter.rs}：一次测量的完整时序}]{../src/meter.rs}

\clearpage
\lstinputlisting[caption={\texttt{src/dsp.rs}：加窗、真有效值与工频估计}]{../src/dsp.rs}

\clearpage
\lstinputlisting[caption={\texttt{src/link.rs}：地址、帧格式与收发}]{../src/link.rs}

\end{document}
